\clearpage
\fancyhead[L]{\small BEER UND MÄDLER.}
\hypertarget{native-657-lower}{}
\begin{center}\small Göttingische gelehrte Anzeigen. Stück 137. Seite 1361..1365. 1841. August 30.\\[3pt]\rule{.72\linewidth}{.3pt}\end{center}

\emph{Beiträge zur physischen Kenntniss der himmlischen Körper im Sonnensysteme von} \textsc{Wilhelm Beer}, \emph{k. preuss. geheimen Commerzienrathe, und Dr.} J. H. \textsc{Mädler}, \emph{kais. russ. Hofrathe.} Weimar 1841. Bei B. F. Voigt. 152 Seiten in Quart, nebst 7 Kupfertafeln.

\emph{Fragments sur les corps célestes du système solaire par} \textsc{Guillaume Beer} \emph{et} \textsc{Jean Henri Mädler.} Paris 1840. Chez Bachelier. 216 Seiten in Quart, mit 7 Tafeln.

Aus denselben Händen, denen wir die classische, vor vier Jahren erschienene, Selenographie verdanken, erhalten wir hier in deutscher und französischer Sprache eine Sammlung kleiner Aufsätze, wovon einige theilweise schon früher bekannt gemacht waren. Sie beziehen sich meistens auf die physische Beschaffenheit des Mondes und der Planeten, und, indem sie das Gepräge sorgfältiger Beobachtung, strenger Prüfung und gesunder Beurtheilung tragen, liefern sie uns manche schätzbare Erweiterung unserer Kenntnisse. Die Gegenstände der einzelnen Aufsätze sind folgende. I. \emph{Ueber die jenseitige Mondhalbkugel.} Die von uns abgekehrte Hälfte der Mondoberfläche zerfällt in zwei Theile, einen, der bei geeigneten Librationsverhältnissen uns sichtbar wird, den andern, der von der Erde aus niemals gesehen werden kann. Die Gestalt dieses zweiten Theils wird hier so beschrieben, er werde von einer Ellipse begrenzt, die eine doppelte Seitenabplattung habe, und seine Grösse wird zu 0.4243 der ganzen Mondoberfläche angegeben. Ref. findet, dass die Gestalt ein \emph{Viereck} ist, und dessen Inhalt 0.4118. Von der physischen Beschaffenheit dieses Theils können wir nun zwar niemals irgend eine positive Kenntniss erhalten: allein die Verff. bemerken mit Recht, dass gar kein Grund vorhanden ist, eine wesentliche Ungleichheit in jener Beziehung zwischen den beiden Hälften des Mondes zu vermuthen, zumal da der durch die Libration sichtbar werdende Theil der jenseitigen Hälfte ganz ähnliche Formationen zeigt, wie die diesseitige. Der einzige Unterschied zwischen den beiden Hälften, den wir mit Gewissheit kennen, ist die Ungleichheit der Beleuchtungsverhältnisse, welche die Verff. mit grosser Ausführlichkeit entwickeln, selbst bis auf die Vortheile für die beobachtende Astronomie, welche die jenseitige Halbkugel des Mondes, wenn der Astronom in Gedanken sich auf dieselbe versetzt, vor anderen Standpunkten voraus haben würde. II. \emph{Ueber die Rillen in der Mondfläche.} Nach \textsc{Schröter}, der zuerst vor 50 \SourceLeaf{658}Jahren die beiden Rillen beim Hyginus und Ariadäus entdeckt hat, sind noch eine Anzahl anderer von \textsc{Pastorf}, \textsc{Gruithuysen} und besonders von \textsc{Lohrmann} wahrgenommen; aber diese ganze Anzahl ist unbedeutend gegen die grosse Menge ähnlicher Formationen, welche die Verff. bei Gelegenheit der Bearbeitung ihrer Mondskarte aufgefunden haben: das hier mitgetheilte Verzeichniss enthält deren 92. III. \emph{Ueber Mondfinsternisse.} Formeln zur Berechnung derselben, besonders in Beziehung auf die Vergrösserung, welche man an dem Halbmesser des Schattenkegels anbringen muss, um die Beobachtungen mit der Rechnung in Uebereinstimmung zu bringen. Die Anwendung auf drei Finsternisse, vom 26. December 1833, 10. Juni 1835 und 13. October 1837 gibt sehr ungleiche Resultate, die erste $\frac{1}{65.4}$, die zweite $\frac{1}{28.3}$, -die dritte $\frac{1}{54.0}$. Als Ursache dieser Erscheinung pflegt man die Atmosphäre der Erde zu betrachten, welche nur einen geringen Theil der Sonnenstrahlen durchgehen lässt und so den Durchmesser des undurchsichtigen Erdkörpers und mithin auch des Schattenkegels vergrössert. Allein diese Erklärungsart ist jedenfalls sehr unbefriedigend, wenn man das Quantitative in der Erscheinung genauer betrachtet. Eine starke Lichtabsorbtion kann nur für die unteren, dichteren Luftschichten eingeräumt werden, während das durch die höheren Luftschichten durchgehende Licht immer weniger geschwächt wird und zugleich vermöge der Refraction den Durchmesser des Schattenkegels vielmehr vermindert. Um die Erscheinung, wie sie in der Mondfinsterniss vom 10. Juni 1835 beobachtet ist, aus jener Ursache zu erklären, müsste man eine undurchsichtige Atmosphäre von 23 Meilen Höhe annehmen. Ohne also jener Ursache eine Mitwirkung abzusprechen, müssen wir zugeben, dass die Beobachter einer Mondfinsterniss das Eintreten eines Punktes der Mondscheibe in den vollen Schatten schon zu sehen glauben, wenn er wirklich noch von einem kleinen Theile der Sonnenscheibe Licht erhält. In dem in Rede stehenden Falle ergibt die Rechnung, dass ohne Berücksichtigung der Schwächung durch die Erdatmosphäre den an der scheinbaren Grenze des Kernschattens liegenden Punkten noch $\frac{1}{43}$ der Sonnenscheibe unbedeckt blieb; unter derselben Einschränkung findet sich $\frac{1}{178}$ für die Finsterniss von 1833, und $\frac{1}{126}$ für die von 1837. Die Aufsätze IV und V enthalten schätzbare Monographien interessanter Mondslandschaften, der Gegend des Kraters \textsc{Schröter} und der Umgegend des Mondnordpols. Der Aufsatz V bezieht sich auf den Saturn und enthält einen sehr verdienstlichen Versuch, die Bahnen der beiden innern Trabanten aus \textsc{Herschel}'s einer neuen Bearbeitung unterzogenen Beobachtungen von 1789 zu bestimmen; eine Entwickelung der Erscheinungen des Saturnringes vom Planeten aus gesehen, und eine Bemerkung über den grauen Aequatorialstreifen auf dem Saturn. Die Verff. deuten auf die Vermuthung hin, dass derselbe vielleicht eine wasserähnliche Flüssigkeit sein könne, die mit der vom Ringe ausgehenden Anziehungskraft in Verbindung stehe. Allein der Behauptung, dass der Ring in den Gegenden, die ihn im Zenith haben, eine viele tausend Mal grössere constante Fluth, als die auf der Erde statt findende, \emph{nothwendig} hervorbringen müsse, können wir nicht beistimmen. In der That hängt das Dasein oder Nichtdasein einer solchen constanten Fluth wesentlich von der Gestalt des festen Theils des Planeten selbst ab, und wenn sich diese, wie wir nach der Analogie für wahrscheinlich halten müssen, unter dem Einflusse aller zu dem Saturnssysteme gehörenden Kräfte, also auch der der Anziehung des Ringes zuzuschreibenden, gebildet hat, so kann eigentlich von einer constanten Fluth gar keine Rede sein. VII. Jupiter; neue Bestimmung seiner Rotationszeit; relative Helligkeit der Trabanten; Messungen des Aequatorial- und Polardurchmessers. VIII. Mars; besonders über die weissen Polarflecken. IX. Venus. Achtzehn Zeichnungen der Gestalt der Lichtgrenze zeigen eigentlich nur, wie schwer es ist, etwas Genaues über die Rotationszeit auszumachen, in Beziehung auf welche die Angaben von \textsc{Cassini} und \textsc{Bianchini} so enorm verschieden sind: die Verff. finden jedoch in jenen Beobachtungen, so wenig sie auch zu einer Entscheidung geeignet sind, eher eine Hindeutung auf die kürzere Rotationszeit. Ueber eine eigenthümliche strahlende Erscheinung, die Hr. \textsc{Mädler} am 7. April 1833 an der Venus bemerkte, lässt sich auf den \SourceLeaf{659}Grund des davon mitgetheilten kein Urtheil fällen. X. Merkur; Messung seines Durchmessers bei dem Vorübergange vor der Sonne 4. Mai 1832. XI. Heliometrische Messungen an einigen Doppelsternen.

\clearpage
\fancyhead[L]{\small BEIMA. DE ANNULO SATURNI.}
\hypertarget{native-659-middle}{}
\begin{center}\small Göttingische gelehrte Anzeigen. Stück 194. Seite 1937..1938. 1842 December 5.\\[3pt]\rule{.72\linewidth}{.3pt}\end{center}

\emph{De annulo Saturni commentatus est} \textsc{Elto Martini Beima}, \emph{mus. hist. nat. publ. conservator.} Leyden 1842. Bei S. und J. Luchtmans. 252 Seiten in Quart, nebst 4 Kupfertafeln.

In dieser Monographie über einen der interessantesten Gegenstände unseres Sonnensystemes ist mit vielem Fleisse. alles gesammelt, was die Geschichte der Entdeckung des Saturnsringes betrifft, und was die Astronomen über seine Lage, Gestalt, Dimensionen, \textperiodcentered{}Masse, physicalische Beschaffenheit, Rotation und Ursprung, so wie über die Kräfte, welche ihn erhalten, und über die Erscheinungen, welche er für seinen Hauptplaneten darbietet, bisher beobachtet, berechnet oder vermuthet haben. Höhere Ansprüche, als auf den Rang einer mit Fleiss gearbeiteten und aus den Quellen geschöpften Compilation kann die Schrift nicht machen, aber dieses Lob kann man ihr beilegen, wenn man gegen eine sehr incorrecte und vernachlässigte Sprache nachsichtig ist, die stellenweise kaum ohne Zuziehung der Quellen, aus welchen übersetzt ist, verständlich wird.
